Los experimentos se realizaron en un equipo con procesador Intel Core i5-12500H,
16 procesadores lógicos y aproximadamente 7,6 GiB de RAM disponibles para Ubuntu
24.04.4 LTS en WSL2. Se utilizó C++17 con g++ 13.3.0 y optimización
\texttt{-O2}. Los datos se generaron con Python 3.12.3 y NumPy 1.26.4;
los gráficos se prepararon con Matplotlib 3.11.1 \cite{matplotlib}.

En ordenamiento se probaron tamaños $10$, $10^3$, $10^5$ y $10^7$, entradas
ascendentes, descendentes y aleatorias, con valores D1 ($0$ a $9$) y D7
($0$ a $10^7$, incluidos). En matrices se probaron tamaños $16$, $64$, $256$
y $1024$, tipos denso, diagonal y disperso, y dominios D0 ($0$ y $1$)
y D10 ($0$ a $9$). Cada configuración tiene tres muestras: a, b y c.
Las matrices dispersas se generan intentando llenar un 10\% de las posiciones;
como pueden repetirse posiciones, la proporción final de no ceros puede ser menor.
Las matrices se almacenan completas, incluso cuando contienen muchos ceros.

Se fijó la semilla 20260911 al iniciar cada generador; en matrices se fijó
tanto en NumPy como en \texttt{random}. Los datos se repiten ejecutando los
scripts completos en el mismo orden y con las mismas versiones. Todos los
resultados de este informe corresponden a la nueva corrida con estas semillas.
Se obtuvieron 288 registros de ordenamiento y 144 de matrices, sin casos
faltantes ni duplicados. Los outputs guardados se contrastaron con un
ordenamiento de referencia y con el producto matricial calculado con NumPy.
Esta comprobación valida los archivos guardados; no valida resultados que
no llegaron a generarse por un timeout.

El tiempo se midió con el reloj \texttt{steady\_clock} de \texttt{std::chrono}, alrededor de la
llamada al algoritmo, sin incluir lectura ni escritura de archivos. La memoria
es el máximo residente del proceso indicado por \texttt{getrusage}, en KiB;
incluye las entradas y no equivale solo a la memoria auxiliar del algoritmo.
Cada ejecución se inicia en un proceso distinto. Las funciones de ordenamiento
retornan el vector por valor, por lo que la medición también puede incluir
el costo de esa interfaz.

Se usaron límites de 30 s para ordenamiento y 60 s para matrices mediante
\texttt{timeout}. Estos límites se aplican al proceso completo, incluida la
entrada/salida, y no son tiempos exactos del algoritmo. Hubo 11 timeouts en
ordenamiento y 18 en matrices. Se presenta la media de las muestras que
terminaron y el rango mínimo--máximo en los gráficos; los timeouts se señalan
aparte y nunca se sustituyen por cero. Hay tres entradas por configuración,
no tres repeticiones de una misma entrada. No se realizaron pruebas de
significancia ni se midieron por separado comparaciones, copias o asignaciones.
